% Dérivé de 2Aw2 (réorganisation des TD de STAT, octobre 2026).
\uuid{Wn8c}
\chapitre{Statistique}
\niveau{L2}
\module{Analyse de données}
\sousChapitre{Autre}
\titre{Approximation normale d'une binomiale et d'une fréquence}
\theme{statistiques, loi binomiale, loi normale}

\auteur{}
\datecreate{2026-10-19}
\organisation{AMSCC}
\difficulte{1}
\contenu{
\texte{$\mathcal N(\mu,\sigma)$ désigne la loi normale de moyenne $\mu$ et d’écart-type $\sigma$.}
\texte{Un restaurant universitaire accueille 600 étudiants et propose risotto ou quiche. On suppose que chacun choisit le risotto avec probabilité $0{,}4$, indépendamment des autres.}
\begin{enumerate}
\item \question{Combien de portions de risotto faut-il prévoir pour que le risque de rupture soit strictement inférieur à $10\%$ ? Utiliser l'approximation de la loi binomiale par une loi normale.}
\indication{Avec $r$ portions, il manque du risotto si $X>r$. Centrer et réduire $X$, puis utiliser le quantile d'ordre $0{,}9$ de la loi normale centrée réduite.}
\reponse{Soit $X$ le nombre d'étudiants souhaitant une portion de risotto. Sous l'hypothèse de choix indépendants, chacun de probabilité $0{,}4$, on a
\[
X\sim\mathcal B(600,0{,}4).
\]
On cherche un entier $r$ tel que $\mathbb P(X>r)<0{,}1$.
L'espérance et la variance de $X$ valent
\[
\mathbb E(X)=600\times0{,}4=240,\qquad
\operatorname{Var}(X)=600\times0{,}4\times0{,}6=144=12^2.
\]
Les quantités $600\times0{,}4=240$ et $600\times0{,}6=360$ étant grandes, on approche la loi de $X$ par $\mathcal N(240,12)$. Ainsi, $\frac{X-240}{12}$ suit approximativement une loi normale centrée réduite, et
\[
\mathbb P(X>r)
=\mathbb P\left(\frac{X-240}{12}>\frac{r-240}{12}\right)
\simeq1-\Phi\left(\frac{r-240}{12}\right).
\]
Dans cette approximation, la condition recherchée s'écrit
\[
\begin{aligned}
1-\Phi\left(\frac{r-240}{12}\right)<0{,}1
&\iff\Phi\left(\frac{r-240}{12}\right)>0{,}9\\
&\iff\frac{r-240}{12}>\Phi^{-1}(0{,}9)\simeq1{,}28155.
\end{aligned}
\]
On obtient donc
\[
r>240+12\times\Phi^{-1}(0{,}9)\simeq255{,}3786.
\]
Le plus petit entier satisfaisant cette condition approchée est $r=256$.
Cette méthode conduit donc à prévoir \textbf{256 portions de risotto}.}
\item \question{On note $F=\frac X{600}$ la proportion d'étudiants qui choisissent le risotto. Donner une loi approchée de $F$, puis $h>0$ tel que $\mathbb P(0{,}4-h\leq F\leq0{,}4+h)\simeq0{,}95$.}
\indication{Diviser $X$ par $600$ divise son espérance et son écart-type par $600$. Pour l'intervalle centré, utiliser le quantile d'ordre $0{,}975$ de la loi normale centrée réduite.}
\reponse{D'après la question précédente, $X$ suit approximativement $\mathcal N(240,12)$. Une transformation affine d'une variable normale reste normale, donc
\[
F=\frac X{600}\approx\mathcal N\left(\frac{240}{600},\frac{12}{600}\right)=\mathcal N(0{,}4\,;\,0{,}02).
\]
On retrouve la forme générale $\mathcal N\left(p,\sqrt{\frac{p(1-p)}n}\right)$ avec $p=0{,}4$ et $n=600$. Alors
\[
\mathbb P(0{,}4-h\leq F\leq0{,}4+h)\simeq2\Phi\left(\frac h{0{,}02}\right)-1=0{,}95
\iff\frac h{0{,}02}=\Phi^{-1}(0{,}975)\simeq1{,}95996,
\]
d'où $h\simeq0{,}0392$. La proportion observée est donc comprise entre environ $36{,}1\%$ et $43{,}9\%$ avec une probabilité proche de $95\%$. Ce calcul prépare les intervalles de confiance d'une proportion, où $p$ est inconnu.}
\end{enumerate}
}
